\section{Discussion}
\label{sec:discussion}

\para{What survives from the hypothesis.}
Three parts hold: NLI is more order-sensitive than sentiment for every model; larger models are usually less order-sensitive on classification; and order variance is flat at higher $K$.
Three parts do not.
The 4-shot threshold disappears once sensitivity is measured with a metric that does not grow with the number of sampled orders; the plateau is reached at $K=2$.
The 15/5-point thresholds fail on average.
And the recency (classification) \vsx primacy (generation) split is reversed in places: recency appears where it was not predicted (summarization with small models) and fails where it was predicted (NLI with mid-size models).

\para{Relation to prior work.}
The recency bias that \citet{zhao2021calibrate} documented in GPT-3 base models appears here only in the smallest instruction-tuned models on \sst; instruction tuning and scale may remove it, although our design cannot separate the two.
\citet{lu2022fantastically} reported that order sensitivity does not vanish with scale.
We find a more mixed picture within one family: it shrinks substantially on easy tasks, persists on \mnli, and grows for summarization at 14B.
The anti-recency effect on NLI is, to our knowledge, not described in prior work and deserves a targeted study.

\para{Practical implications.}
{\bf When is order worth tuning?}
Our reliability analysis gives a direct answer.
For capable models ($\geq$3B) on easy tasks such as \sst, order-induced differences are about 1 point and are no larger than item-level noise, so tuning order on a dev set would overfit.
Where order effects are systematic (\mnli at every size, and summarization at $\geq$3B), averaging over orders or searching for an order is justified.
We also recommend that papers report a pairwise or variance-based sensitivity measure, rather than the best--worst range, when comparing across shot counts or across setups with different numbers of sampled orders, and that they check whether the ``best'' order is reproducible on held-out items.

\para{Error analysis.}
The 0.5B model's summaries usually run to the 80-token cap with repetitive text (median 67 words \vsx 24 for 14B), which explains its low and flat \rougel.
The negative influence on NLI is concentrated in the final position (\figref{fig:labelflip}), which suggests a mechanism specific to the demonstration adjacent to the test item rather than a gradual decay.

\subsection{Limitations}
\label{sec:limitations}

\begin{itemize}[leftmargin=*,itemsep=0pt,topsep=0pt]
    \item {\bf One model family and one template per task.} We test only \qwen (released in 2024) with plain-text prompts and deterministic decoding. Other families, chat-formatted prompts or frontier API models may behave differently; no API models were available on our host.
    \item {\bf Small test sets.} With 200 classification items and 80 dialogues, per-order accuracy carries roughly $\pm$3 points of item-level noise. \Secref{sec:reliability} quantifies this, but larger test sets would sharpen per-cell conclusions.
    \item {\bf Limited orders and sets.} We cap orders at 8 per set, so best--worst ranges are lower bounds on the range over all $K!$ orders ($8!=40{,}320$). 7B and 14B use 10 demonstration sets per $K$ rather than 20, and the 14B probes use only $K=8$ (classification) and $K=6$ (\samsum).
    \item {\bf Narrow generation probe.} The style-marker probe measures one kind of influence: copying a surface marker. We do not test primacy in content selection or in the ordering of facts, and \samsum is our only generation task.
    \item {\bf Non-deterministic inference.} Batched bf16 inference accounts for about a third of the measured \samsum order variance for 0.5B, 1.5B and 7B. The 14B weights ($\approx$30\,GB) did not fit in the free memory of the shared GPUs ($\leq$23\,GB) during the noise-floor session, so the 14B summarization result lacks a noise-floor check.
    \item {\bf Coarse size trends and contamination.} Size trends rest on 5 points. \sst, GLUE and \samsum are public and likely seen in pretraining, which affects absolute accuracy more than order effects.
\end{itemize}

\para{Broader impact.}
This work concerns evaluation methodology. Its main risk is misuse in the opposite direction from our advice: selecting a favorable demonstration order to inflate reported few-shot results. Our reliability test is intended to make such selection visible.
